% Clase 16 --- El conmutador del motor de corriente continua (§2.1).
% "Cada media vuelta hay que dar vuelta el sentido de la corriente... tienen
%  típicamente unas escobillitas para que haga contacto."
% Los dos paneles son la misma posición angular separados por media vuelta,
% vistos a lo largo del eje de giro. Los colores siguen a cada conductor
% (ámbar = 1, azul = 2) para que se vea que rotaron, y los símbolos (x) / (.)
% muestran que la corriente en cada conductor se invirtió. Moraleja: en una
% posición dada la corriente es siempre la misma, así que el par nunca cambia
% de sentido y la espira gira en lugar de oscilar.
% El sub-dibujo del conmutador va con \providecommand y ANTES del tikzpicture
% (si se define dentro de un scope no se ve desde el otro panel).
\providecommand{\conmutador}[1]{%
  \draw[figamber, line width=2.2pt] ({#1+7}:0.62) arc ({#1+7}:{#1+173}:0.62);
  \draw[figblue, line width=2.2pt] ({#1+187}:0.62) arc ({#1+187}:{#1+353}:0.62);
  \node[figlblaux, figamber] at ({#1+90}:0.95) {$1$};
  \node[figlblaux, figblue] at ({#1+270}:0.95) {$2$};
  \draw[figgray, fill=figgray!25, line width=0.6pt]
    (-0.95,-0.15) rectangle (-0.70,0.15);
  \draw[figgray, fill=figgray!25, line width=0.6pt]
    (0.70,-0.15) rectangle (0.95,0.15);
}
\begin{center}
\begin{tikzpicture}[scale=0.95]

  % =============== (a) ===============
  \begin{scope}
    \draw[figvecaux, figgray] (-2.9,2.8) -- (2.4,2.8);
    \node[figlblaux, anchor=south east] at (2.4,2.85) {$\vec B$};

    % traza del plano de la espira
    \draw[figaux] (-1.5,-1.5) -- (1.5,1.5);

    % conductor 1 (ámbar): saliente;  conductor 2 (azul): entrante
    \node[figamber, scale=1.15] at (1.34,1.34) {$\odot$};
    \node[figlbl, figamber, anchor=south west] at (1.42,1.40) {$1$};
    \node[figblue, scale=1.15] at (-1.34,-1.34) {$\otimes$};
    \node[figlbl, figblue, anchor=north east] at (-1.42,-1.40) {$2$};

    % fuerzas: saliente -> hacia arriba;  entrante -> hacia abajo
    \draw[figvecres, line width=1.2pt] (1.34,1.62) -- (1.34,2.45);
    \node[figlbl, figred, anchor=east] at (1.24,2.05) {$\vec F$};
    \draw[figvecres, line width=1.2pt] (-1.34,-1.62) -- (-1.34,-2.45);
    \node[figlbl, figred, anchor=west] at (-1.24,-2.05) {$\vec F$};

    % sentido de giro
    \draw[figvecaux, figred, line width=1pt,
          -{Latex[length=2.2mm,width=1.6mm]}] (115:2.25) arc (115:162:2.25);
    \node[figlblaux, figred, anchor=south east] at (-1.35,1.75) {gira};

    \conmutador{45}
    \node[figlblaux, anchor=north, align=center] at (0.45,-0.95)
      {escobillas fijas};

    % batería
    \draw[figgray, line width=0.8pt]
      (-0.95,0) -- (-2.65,0) -- (-2.65,-3.2) -- (-0.45,-3.2);
    \draw[figgray, line width=0.8pt]
      (0.95,0) -- (2.65,0) -- (2.65,-3.2) -- (0.45,-3.2);
    \draw[figgray, line width=1.5pt] (-0.16,-3.45) -- (-0.16,-2.95);
    \draw[figgray, line width=0.9pt] (0.16,-3.32) -- (0.16,-3.08);
    \draw[figgray, line width=0.8pt] (-0.45,-3.2) -- (-0.16,-3.2);
    \draw[figgray, line width=0.8pt] (0.16,-3.2) -- (0.45,-3.2);
    \node[figlblaux, anchor=north] at (0,-3.55) {$\varepsilon$};

    \node[figlbl, anchor=north, align=center] at (0,-4.15)
      {(a) antes de la media vuelta};
  \end{scope}

  % =============== (b) ===============
  \begin{scope}[xshift=7.6cm]
    \draw[figvecaux, figgray] (-2.9,2.8) -- (2.4,2.8);
    \node[figlblaux, anchor=south east] at (2.4,2.85) {$\vec B$};

    \draw[figaux] (-1.5,-1.5) -- (1.5,1.5);

    % ahora arriba está el 2, y por él sale la corriente
    \node[figblue, scale=1.15] at (1.34,1.34) {$\odot$};
    \node[figlbl, figblue, anchor=south west] at (1.42,1.40) {$2$};
    \node[figamber, scale=1.15] at (-1.34,-1.34) {$\otimes$};
    \node[figlbl, figamber, anchor=north east] at (-1.42,-1.40) {$1$};

    \draw[figvecres, line width=1.2pt] (1.34,1.62) -- (1.34,2.45);
    \node[figlbl, figred, anchor=east] at (1.24,2.05) {$\vec F$};
    \draw[figvecres, line width=1.2pt] (-1.34,-1.62) -- (-1.34,-2.45);
    \node[figlbl, figred, anchor=west] at (-1.24,-2.05) {$\vec F$};

    \draw[figvecaux, figred, line width=1pt,
          -{Latex[length=2.2mm,width=1.6mm]}] (115:2.25) arc (115:162:2.25);
    \node[figlblaux, figred, anchor=south east, align=right] at (-1.35,1.75)
      {sigue\\[-0.25em] girando};

    \conmutador{225}
    \node[figlblaux, anchor=north, align=center] at (0.5,-0.95)
      {cada anillo cambió\\[-0.25em] de escobilla};

    \draw[figgray, line width=0.8pt]
      (-0.95,0) -- (-2.65,0) -- (-2.65,-3.2) -- (-0.45,-3.2);
    \draw[figgray, line width=0.8pt]
      (0.95,0) -- (2.65,0) -- (2.65,-3.2) -- (0.45,-3.2);
    \draw[figgray, line width=1.5pt] (-0.16,-3.45) -- (-0.16,-2.95);
    \draw[figgray, line width=0.9pt] (0.16,-3.32) -- (0.16,-3.08);
    \draw[figgray, line width=0.8pt] (-0.45,-3.2) -- (-0.16,-3.2);
    \draw[figgray, line width=0.8pt] (0.16,-3.2) -- (0.45,-3.2);
    \node[figlblaux, anchor=north] at (0,-3.55) {$\varepsilon$};

    \node[figlbl, anchor=north, align=center] at (0,-4.15)
      {(b) media vuelta después};
  \end{scope}

\end{tikzpicture}\\[0.5em]
{\small\itshape Sin conmutador el motor sería ``aburrido'': pasada la posición
de alineación el par se invierte y la espira oscila en vez de girar. El truco es
que la corriente que entrega la batería \textbf{no} cambia de signo ---eso sería
difícil---; lo que cambia es la conexión. Los dos medios anillos giran
solidarios con la espira y, al pasar por la posición crítica, cada uno queda
apoyado en la escobilla opuesta. Comparando los paneles: el conductor $1$ pasó
de arriba a abajo y a la vez su corriente se invirtió de saliente a entrante, de
modo que \emph{en cada posición} la corriente ---y con ella la fuerza--- es
siempre la misma, y el par mantiene el sentido vuelta tras vuelta. El motor de
corriente alterna funciona igual, sólo que ahí la inversión la trae la propia
corriente, que oscila sola. Como $\tau\propto\sin\theta$ no es constante a lo
largo de la vuelta, en la práctica se agregan dispositivos que estabilicen el
giro para que no salga a tirones.}
\end{center}
